\documentclass[../main]{subfiles}
\begin{document}
\chapter*{Problemas Resueltos}
\section{TP3:Principio 1}
\begin{enumerate}
 \item .

 \info{Datos}{
 
 $A=0,00234\dfrac{kcal}{kgm}$
 
 $W= 2562 kgm$
 
 }
 
 $Q = U_2 – U_1 + A W$
 
Si $\Delta U$ es constante quiere decir que $U_2 = U_1$ y entonces $U_2 – U_1 = 0$

Por lo tanto $Q = A \times W$

 $Q = + 0,00234\dfrac{kcal}{kgm} \times 2562 kgm = 6 Kcal$

\item Si la energía interna del gas es constante, entonces por la Primera Ley de la Termodinámica:

$\Delta U = Q - W = 0$

donde $\Delta U$ es el cambio en la energía interna del gas, $Q$ es la cantidad de calor entregada al gas y $W$ es el trabajo realizado sobre el gas.

En este caso, se nos da $Q = 30\ \text{kcal}$. Para expresar la cantidad de calor en unidades de energía más pequeñas, como los joules, podemos utilizar la relación:

$1\ \text{kcal} = 4184\ \text{J}$

Entonces, $Q = 30\ \text{kcal} \cdot 4184\ \text{J/kcal} = 125,520\ \text{J}$

Sustituyendo en la ecuación de la Primera Ley de la Termodinámica, obtenemos:

$W = Q = 125,520\ \text{J}$

Por lo tanto, el trabajo externo que es necesario ejercer sobre el gas cuando se comprime y entrega $30\ \text{kcal}$ de energía interna constante es de aproximadamente $125,520\ \text{J}$.
 \item El trabajo realizado sobre el gas al comprimirlo es:

$W = Fd = pA d = 2\ \text{atm} \cdot 101325\ \text{Pa/atm} \cdot 0.01\ \text{m}^2 \cdot 0.3\ \text{m} = 6079.5\ \text{J}$

La cantidad de calor transferida al gas es $Q = 5\ \text{kcal} \cdot 4184\ \text{J/kcal} = 20,920\ \text{J}$

Por lo tanto, la variación de la energía interna del gas es:

$\Delta U = Q - W = 20,920\ \text{J} - 6079.5\ \text{J} = 14,840.5\ \text{J}$

La respuesta en kilocalorías se obtiene dividiendo el resultado en julios entre $4184\ \text{J/kcal}$:

$\Delta U = 14,840.5\ \text{J} / 4184\ \text{J/kcal} = 3.54\ \text{kcal}$

Por lo tanto, la variación de la energía interna del gas cuando es comprimido es de aproximadamente $3.54\ \text{kcal}$.

\item .
\info{Problema 4}{

$p = 4 atm$

$\Delta V = 0,5 m^3$

$Q = 0 kcal$

$A= 0,00234 \frac{kcal}{kgm}$
}
Como se trata de una transformación abierta se aplica la expresión:

 $\Delta U = U_2 – U_1 = Q – A W$

En este caso $Q = 0$ luego nos queda: $\Delta U = – A W$
Tenemos un trabajo de expansión 1-2 positivo y representado por el área: 1-2-3-4 y su valor es:

$ W = p ( V_2 – V_1 )$ pero aquí el $\Delta V$ es dato = $0,5 m^3$ 
 
Quedando entonces: 
$W = p \Delta V = 4 atm \times 0,5 m^3 = 40000 \times 0,5 m^3 = \fbox{20000 kgm}$
 

%problema 5
\item 
Datos:

$W_{1-3} = 50000 Kgm$

$W_{1-2-3} = 10000 Kgm$

$Q_{1-3} = 800 kcal$

Para la transformación abierta $1-3$ aplicamos:

 $Q_{1-3} = \Delta U_{1-3} + A W_{1-3}$

Aquí tenemos como datos el valor de Q y del trabajo por lo que el valor de $\Delta U$ será:

 $\Delta U_{1-3} = Q1-3 – A W_{1-3}$
 
 $\Delta U_{1-3} = 800 Kcal – 0,00234\frac{Kcal}{kgm} \times 50000 kgm$
 
 $\Delta U_{1-3} = 800 Kcal – 117 Kcal = 683 Kcal$

Como nos da un valor positivo significa que el gas también aumenta su temperatura.

Para la transformación abierta 1- 2 - 3 también aplicamos: 

$Q_{1-2-3} =\Delta U_{1-2-3} + A W_{1-2-3}$

Pero aquí no tenemos como dato el valor de $Q_{1-2-3}$ solo el del trabajo $W_{1-2-3}$

Para poder resolver sabemos que la energía interna en un punto es única por lo tanto a diferencia del trabajo por cualquiera que sea el camino es $\Delta U_{1-3} = \Delta U_{1-2-3}$
 
 $Q_{1-3} = \Delta U1-3 + A W1-3$

 $Q_{1-3} = 683 Kcal + 0,00234 \times 10000 kgm$
 
 $Q_{1-3} = 683 Kcal + 23,4 Kcal = 706,4 Kcal$

Como nos da un valor positivo significa que el gas también aumenta su temperatura 

\item Para calcular la diferencia de entalpía por unidad de peso que experimentó la masa de aire, podemos utilizar la ecuación:

$\Delta h = \Delta u + \Delta (p/\gamma)$

donde $\Delta h$ es la diferencia de entalpía, $\Delta u$ es la diferencia de energía interna, $\Delta (p/\gamma)$ es la diferencia de presión específica y se define como $(p_e/\gamma_e) - (p_i/\gamma_i)$.

Primero, calculemos la diferencia de presión específica:

$\Delta (p/\gamma) = (p_e/\gamma_e) - (p_i/\gamma_i) =$

$ (5 atm / 4 kg/m^3) - (1 atm / 1.25 kg/m^3) =$

$\dfrac{5 atm \times (101325\frac{Pa}{atm})} { 4 kg/m^3} - \dfrac{1 atm \times 101325\frac{Pa}{atm}}{ 1.25 kg/m^3} =126656.25 J/kg - 81060J/kg=$

$45596.25J/kg \times \dfrac{1kcal}{4184J}=10,9kcal/kg$

Luego, podemos calcular la diferencia de entalpía:

$\Delta h= \Delta u + \Delta p/\gamma$

$\Delta h=32kcal/kg-4kcal/kg + 10,9 kcal/kg =$

$\Delta h=28kcal/kg + 10,9kcal/kg=\fbox{38,9kcal/kg}$

\info{Entalpía}{La entalpía es una propiedad termodinámica que describe la energía total de un sistema termodinámico, incluyendo tanto su energía interna como la energía que el sistema intercambia con su entorno en forma de trabajo y calor.

La entalpía se representa por la letra H y se define como H = U + PV, donde U es la energía interna del sistema, P es la presión del sistema y V es su volumen. La entalpía es una función de estado, lo que significa que su valor final depende solo del estado inicial y final del sistema y no del camino seguido para llegar de uno a otro.}


 
\end{enumerate}

\end{document}